% Synthetic mini manuscript for trying the number guard.
% All numbers are invented. Frozen run: 2025-03-14_101500
\documentclass[11pt]{article}
\usepackage[margin=2.5cm]{geometry}
% \res{label}{value}: bound to one line of the frozen run (label-level binding)
% \ext{key}{value}:   bound to one row of external_numbers.csv (exact)
\newcommand{\res}[2]{#2}
\newcommand{\ext}[2]{#2}

\title{Remote Work and Commuting: A Synthetic Example}
\author{Example Author}
\date{}

\begin{document}
\maketitle

\section{Data}\label{sec:data}

The analysis sample comprises \res{n_persons}{4,812} persons and
19,230 person-waves from two waves of a synthetic household panel
(frozen run 2025-03-14\_101500). The share of employees working remotely at
least one day a week rose from 22.4\% in the first wave to 28.7\% in the
second, an increase of \res{diff_share_w2_w1}{6.3} percentage points. An
external benchmark puts the national share at \ext{benchmark_remote_share}{24.0}\%
in 2023.

\section{Results}\label{sec:results}

In the fixed-effects model (Table~\ref{tab:fe}), starting remote work is
associated with a decrease in weekly commuting hours of
\res{coef_remote}{1.84} (SE 0.41; $N = 19{,}230$).

\begin{table}[h]
\centering
\caption{Fixed-effects estimate (synthetic).}\label{tab:fe}
\begin{tabular}{lr}
\hline
Remote work & $-1.84$ \\
 & $(0.41)$ \\
Person-waves & 19,230 \\
\hline
\end{tabular}
\end{table}

\end{document}
